\begin{pastexam}[label={px:ntu-110-esoe-ds-09}]{臺大工科海洋所 110 學年度（A）第 9 題}
（追蹤）

\textbf{(10 \%)} For the inputs $\{4,11,17,14,23,8,31\}$ and a hash function $h(x)=x\ \bmod\ 7$, show
the linear probing hash table of size 7.
\end{pastexam}

\begin{pastsoln}
照題目列出的順序插入。 線性探測：從 $h(x)$ 開始往後一格一格找空位，超過第 $6$ 格繞回第 $0$ 格
（資結書第~12.3~節）。
\begin{center}
\begin{tabular}{ccl}
\toprule
$x$ & $h(x) = x \bmod 7$ & 探測的格子 $\to$ 放在 \\
\midrule
4 & 4 & $4 \to 4$ \\
11 & 4 & $4$（4）, $5 \to 5$ \\
17 & 3 & $3 \to 3$ \\
14 & 0 & $0 \to 0$ \\
23 & 2 & $2 \to 2$ \\
8 & 1 & $1 \to 1$ \\
31 & 3 & $3$（17）, $4$（4）, $5$（11）, $6 \to 6$ \\
\bottomrule
\end{tabular}
\end{center}
括號裡是那一格已經放的鍵。 最後的表（剛好 $7$ 個鍵填滿 $7$ 格）：
\begin{center}
\begin{tabular}{*{7}{c}}
\toprule
0 & 1 & 2 & 3 & 4 & 5 & 6 \\
\midrule
14 & 8 & 23 & 17 & 4 & 11 & 31 \\
\bottomrule
\end{tabular}
\end{center}
碰撞發生在 $11$（和 $4$ 同在第 $4$ 格）和 $31$（和 $17$ 同在第 $3$ 格，又一路被 $4$、$11$ 擋住）。 $31$ 探測
了 $4$ 次，這就是線性探測的「群聚」：第 $3$--$5$ 格已經連成一段，落在這段裡的鍵都要走到段尾。 表滿
了以後再插入任何鍵都會繞一圈找不到空位，實際使用時要在載入因子太高之前擴大表。
\end{pastsoln}
